\begin{abstract}
\noindent
模型自动扩缩容是实现无服务器模型即服务的关键机制，
但它在扩缩速度与用于缓存参数的存储/内存占用之间存在根本性权衡，
因而难以满足跨多主机、频繁发生的扩容需求。
核心问题在于数据平面过慢：
在参数加载完成之前，新扩出的实例始终处于停机状态。

本文首先指出，如果通过 GPU 之间的计算网络加载参数，
那么即使完全不依赖缓存或仅使用 $O(1)$ 缓存，
也能让数据平面变得\emph{快速}，原因在于：
(1) 该路径的速度与主机缓存相当，却长期未被充分利用；
(2) 借助针对网络优化的多播，扩多个实例时仍只需要零或 $O(1)$ 缓存。
其次，我们通过将推理扩缩容的抽象从粗粒度的实例级下沉到细粒度的层级，
使自动扩缩容变得\emph{在线}。
这样一来，无需等待参数全部加载完成，
就可以把过载服务实例中的部分层计算卸载到新扩出的实例上。

在真实工作负载下，
我们的系统 {\sys} 相比最先进的自动扩缩容系统 ServerlessLLM，
最多可将尾延迟降低 94\,\%；
而与不支持自动扩缩容、但满足相同服务等级协议的 DistServe 与 vLLM 相比，
{\sys} 可将服务所消耗的 GPU 时间降低 49\,\%。
\end{abstract}
